%%%%%%%%%%%%%%%%%%%%%%%%%%%%%%%%%%%%%%%%%%%%%%%%%%%%%%%%%%%%%%%%%
% Project     : LaTeX doc Vorlage für Windows ProTeXt mit TexMakerX
% File        : PA.tex Rev. 00 (Hauptfile)
%%%%%%%%%%%%%%%%%%%%%%%%%%%%%%%%%%%%%%%%%%%%%%%%%%%%%%%%%%%%%%%%%

\documentclass[oneside,12pt,parskip=half,american,numbers=noenddot]{scrreprt}

\usepackage[utf8]{inputenc}
\usepackage[american]{babel}
\usepackage{graphicx}
\usepackage{url}
\usepackage{geometry}
\usepackage{amsmath,amssymb}
\usepackage{longtable}
\usepackage{float}
\usepackage{setspace}
\usepackage{tabularx}
\usepackage{makecell}
\usepackage{newtxtext,newtxmath}

\usepackage{tikz}
\usetikzlibrary{shapes.geometric, arrows.meta, positioning}
\usepackage[plainpages=false]{hyperref}

% Bibliography
\usepackage[style=apa, backend=biber]{biblatex}
\addbibresource{references.bib}
\addbibresource{manuell.bib}

% Page layout (YOUR original geometry)
\geometry{
  a4paper,
  top=1.5cm,
  bottom=1.5cm,
  left=3cm,
  right=3cm,
  headheight=14pt,
  headsep=0.6cm,
  footskip=1.25cm,
  includeheadfoot
}

% Header/footer vertical placement (tune to taste)
\setlength{\headsep}{7mm}     % smaller -> header closer to top edge
\setlength{\footskip}{10mm}   % smaller -> footer closer to bottom edge

% Format list of figures/tables
\usepackage{tocloft}
\renewcommand{\cftfigpresnum}{Figure~}
\renewcommand{\cfttabpresnum}{Table~}
\renewcommand{\cftfigaftersnum}{\hspace{0.5em}}
\renewcommand{\cfttabaftersnum}{\hspace{0.5em}}
\settowidth{\cftfignumwidth}{Figure~99.99}
\settowidth{\cfttabnumwidth}{Table~99.99}

%%%%%%%%%%%%%%%%%%%%%%%%%%%%%%%%%%%%%%%%%%%%%%%%%%%%%%%%%%%%%%%%%
% === KOMA header/footer setup (scrlayer-scrpage ONLY) ===
\usepackage[automark]{scrlayer-scrpage}

% --- FRONT MATTER: no lines, page number only (must work on chapter opening pages too)
\newpairofpagestyles{frontmatter}{
  \clearpairofpagestyles
  \ihead{} \ohead{}
  \ifoot{} \ofoot{\pagemark}
}

% --- MAIN MATTER: lines + chapter title in header + page number in footer
\newpairofpagestyles{mainmatter}{
  \clearpairofpagestyles
  \automark[chapter]{chapter}
  \ihead{\leftmark}
  \ohead{}
  \ifoot{} \ofoot{\pagemark}
}

% Make the chapter-opening pages behave like the chosen pagestyle
% (This is THE key for ToC/LoF/LoT page 1 and for chapter page 1.)
\renewcommand*{\chapterpagestyle}{frontmatter} % will be switched later to mainmatter

% Chapter spacing (reduce the big drop everywhere)
\RedeclareSectionCommand[
  font=\huge,
  beforeskip=0pt plus 1pt minus 1pt,
  afterskip=8pt plus 1pt minus 1pt
]{chapter}

%%%%%%%%%%%%%%%%%%%%%%%%%%%%%%%%%%%%%%%%%%%%%%%%%%%%%%%%%%%%%%%%%
% --- Begin document ---
\begin{document}


%%%%%%%%%%%%%%%%%%%%%%%%%%%%%%%%%%%%%%%%%%%%%%%%%%%%%%%%%%%%%%%%%
% --- Title page (no numbering) ---
\begin{titlepage}
    \thispagestyle{empty} % No page number here
    \centering
    \vspace*{1.8cm} % pushes everything down a bit from top

    % --- University Header ---
    {\large
    \textbf{Zurich University of Applied Sciences} \\
    School of Management and Law \\
    \vspace{0.4cm}
    Department Banking, Finance, Insurance \\
    }

    \vfill

    % --- Paper Title Block ---
    {\Large \textbf{Deep Learning FS 26}} \\[1.0cm]

    {\LARGE\bfseries
    Fraud Detection
    }

    \vfill

    % --- Author and Supervisor Information ---
    {\large
    \textbf{Submitted by} \\[0.3cm]
    Francis Du \\
    Max Heinzer \\
    Dominik Kaufmann \\[1.2cm]

    \textbf{Supervised by} \\[0.3cm]
    Bledar Fazlija \\[1.4cm]

    Winterthur, May 12\textsuperscript{th}, 2026
    }

    \vspace*{1.2cm} % ensures bottom spacing adapts gracefully
\end{titlepage}

%%%%%%%%%%%%%%%%%%%%%%%%%%%%%%%%%%%%%%%%%%%%%%%%%%%%%%%%%%%%%%%%%
% --- FRONT MATTER: Roman numbering, NO lines, page numbers visible ---
\pagenumbering{Roman}
\setcounter{page}{1}

\pagestyle{frontmatter}
\renewcommand*{\chapterpagestyle}{frontmatter}
\KOMAoptions{headsepline=false, footsepline=false}

\onehalfspacing

% === Management Summary ===
\chapter*{Management Summary}
\thispagestyle{frontmatter}
\addcontentsline{toc}{chapter}{\textbf{Management Summary}}

This space is reserved for the management summary.

\newpage

% === Table of Contents ===
\renewcommand{\contentsname}{\textbf{Table of Contents}}
\addcontentsline{toc}{chapter}{\textbf{Table of Contents}}
\tableofcontents
\thispagestyle{frontmatter}
\newpage

% === List of Tables ===
\renewcommand{\listtablename}{\textbf{List of Tables}}
\addcontentsline{toc}{chapter}{\textbf{List of Tables}}
\listoftables
\thispagestyle{frontmatter}
\newpage

% === List of Figures ===
\renewcommand{\listfigurename}{\textbf{List of Figures}}
\addcontentsline{toc}{chapter}{\textbf{List of Figures}}
\listoffigures
\thispagestyle{frontmatter}
\newpage

%%%%%%%%%%%%%%%%%%%%%%%%%%%%%%%%%%%%%%%%%%%%%%%%%%%%%%%%%%%%%%%%%
% --- MAIN MATTER: Arabic numbering, WITH lines (also on chapter page 1) ---
\pagenumbering{arabic}
\setcounter{page}{1}

\pagestyle{mainmatter}
\renewcommand*{\chapterpagestyle}{mainmatter}
\KOMAoptions{headsepline=true, footsepline=true}

\onehalfspacing

% === Main content ===
\chapter{Introduction}\label{Introduction}

\section{Background}\label{Background}
Credit card fraud detection is a critical problem in modern financial systems, characterized by significant consequences and challenging data structures. The rarity of fraudulent transactions relative to legitimate ones introduces severe class imbalance, complicating both model training and evaluation. A model that appears highly accurate can still be practically useless if it misses most fraud cases or overwhelms analysts with false alarms. \textcite{gunjan_credit_2022} describe this imbalance as a central obstacle in credit card fraud detection, since naive classifiers can achieve deceptively strong aggregate metrics by predicting the majority class almost everywhere.

As digital payments have expanded, fraud has increasingly shifted toward card-not-present transactions, where authorization relies primarily on card details and contextual signals rather than physical possession of the card \parencite{xuan_random_2018}. This raises both the economic and operational stakes of effective fraud detection. At the same time, the problem is not static: fraudsters continuously adapt their strategies, which means that models must identify rare events under changing behavioral patterns rather than under a fixed and stable data-generating process \parencite{afriyie_supervised_2023}. Because our dataset provides transaction-level fraud labels, a supervised learning framework remains a natural starting point for comparing predictive approaches and evaluating imbalance-handling techniques.

In our previous work, we examined the same dataset using classical supervised machine-learning models under a comparable evaluation design. Specifically, we evaluated Random Forest and Support Vector Machine classifiers on both the original imbalanced data and a SMOTE-resampled version of the training data. Although the Random Forest achieved the strongest overall balance among the earlier models, performance remained limited, and SMOTE did not improve generalization to the temporally separated test set. In some cases, oversampling substantially increased recall only at the cost of a dramatic loss in precision, producing an operationally unrealistic number of false positives. These findings suggest that class imbalance cannot be addressed as a purely mechanical preprocessing problem. Instead, the central issue is whether a different model class can learn more informative structure from the same transactions without recreating the same false-positive burden.

Recent literature suggests that deep learning is a promising next step in this context. \textcite{mienye_deep_2024} note that, although most credit card fraud studies have historically relied on traditional machine learning, deep learning applications have increased substantially in recent years. In particular, recurrent architectures such as RNNs, LSTMs, and GRUs are especially relevant for fraud detection because they are designed to model sequential data and temporal relationships in transaction streams \parencite{mienye_deep_2024}. Similarly, \textcite{zhou_deep_2025} argue that fraud detection benefits from models that can capture both short-term fluctuations and longer-term dependencies in transaction behavior. This is particularly relevant in our setting, where fraudulent activity may emerge not only through isolated feature values but also through changes in timing, frequency, and behavioral sequence.

At the same time, the recent deep learning literature makes clear that class imbalance remains a major challenge even when more expressive models are used. \textcite{mienye_deep_2024} emphasize that resampling, cost-sensitive learning, and loss-function adaptation remain important strategies in deep learning-based fraud detection, while \textcite{sekhar_nama_credit_2026} likewise stress that fraud datasets remain severely imbalanced and that sampling choices can materially influence downstream performance. This is directly relevant for the present paper, since one of the main lessons from our earlier machine learning study was that SMOTE did not help in our original setting and, in some cases, substantially worsened generalization.

Building on these results, the present paper shifts from classical machine-learning models to a deep learning approach based on Long Short-Term Memory (LSTM) networks. Using the same dataset and a comparable evaluation logic, we examine whether a sequence model can better capture temporal and behavioural dependencies in transaction data and thereby improve fraud detection on unseen observations. At the same time, we revisit SMOTE in this new modelling context. Since oversampling performed poorly in our earlier study, a key question is whether synthetic balancing becomes more useful when combined with a model that is explicitly designed to process sequential information, or whether the same generalization problem persists.

\section{Research Question}\label{Research Question}
This study examines whether a deep learning approach can improve credit card fraud detection under severe class imbalance. Using the same dataset and a comparable evaluation design as in our previous machine-learning study, we evaluate an LSTM model trained on the original imbalanced data and on training data balanced using SMOTE.

The primary research question is:

\textit{Can an LSTM-based model improve fraud detection on this dataset relative to our earlier classical benchmarks, and does applying SMOTE to the training data improve or worsen out-of-sample performance compared with training on the original imbalanced data?}

Because fraud detection is an operational classification problem, the answer is assessed primarily through fraud recall, fraud precision, and the broader precision--recall trade-off on the temporally separated test set.

\pagebreak
\chapter{Methods}

Figure~\ref{fig:ml_pipeline} provides an overview of the modelling pipeline adopted in this study. The workflow is designed to maintain a strict separation between model development and final evaluation, thereby reducing information leakage and overly optimistic performance estimates. Following exploratory data analysis, feature engineering, sequence construction, and preprocessing, the development dataset (\texttt{fraudTrain.csv}) is split into training and validation subsets. Model fitting, hyperparameter tuning, and interim evaluation are conducted exclusively within this development sample.

In contrast to our previous study, which focused on Random Forest and Support Vector Machine classifiers, the present paper evaluates an LSTM-based deep learning model under two training conditions: the original imbalanced data and a SMOTE-resampled training set. The held-out file \texttt{fraudTest.csv}, which is temporally separated from the development data, is reserved exclusively for final out-of-sample evaluation.

This design allows us to assess not only whether an LSTM improves fraud detection relative to our earlier benchmarks, but also whether SMOTE contributes to or harms generalization in the deep-learning setting.

\begin{figure}[H]
    \centering
    \includegraphics[width=1\textwidth]{images/graphs/ml_pipeline.pdf}
    \caption{ML Pipeline}
    \label{fig:ml_pipeline}
\end{figure}

\section{Dataset Description}
\label{sec:dataset_description}

The analysis is conducted on a publicly available credit card transaction dataset obtained from Kaggle \parencite{kaggle_fraud_detection_dataset}. The dataset is provided as two separate files, \texttt{fraudTrain.csv} and \texttt{fraudTest.csv}, which are temporally disjoint and share an identical feature schema.  Together, the two files correspond to an approximate 70/30 split between development and test data.

The file \texttt{fraudTrain.csv} contains 1,296,675 observations and is used as the development dataset. This dataset is further split into training and validation subsets using an 80/20 stratified split, as described in Section~\ref{sec:data_split}, and is used for model fitting, preprocessing, and hyperparameter tuning. The file \texttt{fraudTest.csv} contains 555,719 observations and serves as a held-out test set. It is not used during model development and is reserved exclusively for final out-of-sample evaluation.

\subsubsection{Training Set}
The training dataset contains approximately 1.3 million individual credit card transactions recorded between early 2019 and mid-2020. Each observation corresponds to a single transaction and is associated with a binary target variable, \texttt{is\_fraud}, indicating whether the transaction was fraudulent. The dataset exhibits a pronounced class imbalance, with fraudulent transactions accounting for a small fraction of all observations.

The data contain a heterogeneous set of transactional, temporal, geographical, and customer-level variables. Transaction attributes include the amount (\texttt{amt}), merchant identifier (\texttt{merchant}), merchant category (\texttt{category}), and transaction timestamps (\texttt{trans\_date\_trans\_time}, \texttt{unix\_time}). Additional variables describe customer demographics and location, including gender, city, state, ZIP code, city population, occupation, and date of birth. Geographic coordinates are available for both the cardholder and the merchant, enabling the construction of spatially informed features. At the same time, the dataset also includes several identifier-like fields, such as credit card numbers, transaction identifiers, names, and street-level addresses, which require careful handling to avoid memorization and leakage. Legitimate transactions account for approximately 99\% of all observations, while fraudulent transactions represent only a very small minority, underscoring the extreme imbalance characteristic of real-world fraud detection tasks.


\subsubsection{Test Set}
The held-out test dataset contains approximately 550,000 transactions recorded from mid-2020 to early 2021. It follows the same feature structure as the training data, including the same set of transactional, categorical, temporal, and geographical variables. The distribution of key attributes, such as merchant categories, transaction amounts, and demographic characteristics, is broadly consistent with the training set, while reflecting natural temporal variation. The test set is used exclusively for final model evaluation and is not involved in feature engineering, encoding, or model training.

\subsubsection{Feature Overview.}
An overview of the variables available in both datasets is provided below:

\begin{table}[H]
\centering
\caption{Overview of features in the credit card fraud dataset.}
\label{tab:dataset_features}
\small
\renewcommand{\arraystretch}{1.2}
\begin{tabular}{ll}
\hline
\textbf{Feature} & \textbf{Description} \\
\hline
\texttt{index} & Unique identifier for each transaction \\
\texttt{trans\_date\_trans\_time} & Transaction date and time \\
\texttt{cc\_num} & Credit card number of the customer \\
\texttt{merchant} & Merchant name \\
\texttt{category} & Merchant category \\
\texttt{amt} & Transaction amount \\
\texttt{first}, \texttt{last} & Cardholder first and last name \\
\texttt{gender} & Cardholder gender \\
\texttt{street}, \texttt{city}, \texttt{state}, \texttt{zip} & Cardholder address information \\
\texttt{lat}, \texttt{long} & Cardholder geographic coordinates \\
\texttt{city\_pop} & Population of the cardholder’s city \\
\texttt{job} & Cardholder occupation \\
\texttt{dob} & Cardholder date of birth \\
\texttt{trans\_num} & Unique transaction identifier \\
\texttt{unix\_time} & Unix timestamp of the transaction \\
\texttt{merch\_lat}, \texttt{merch\_long} & Merchant geographic coordinates \\
\texttt{is\_fraud} & Fraud indicator (target variable) \\
\hline
\end{tabular}
\end{table}

The presence of high-cardinality identifiers and sensitive personal attributes motivates careful feature selection and leakage prevention strategies, which are discussed in subsequent sections.


\section{Exploratory Data Analysis}
\label{sec:eda}

This section describes the exploratory analysis conducted to understand the structure of the data, class imbalance, and key characteristics of transactional, temporal, and categorical variables.

\subsection{Class Imbalance}
The dataset exhibits extreme class imbalance, with fraudulent transactions representing well below one percent of all observations. This imbalance makes accuracy-based evaluation misleading and motivates the use of precision--recall metrics, stratified sampling, and explicit imbalance-handling strategies throughout model development. Figure~\ref{fig:eda_class_imbalance} illustrates the extreme rarity of fraudulent transactions, which complicates reliable fraud detection and performance assessment.

\begin{figure}[H]
    \centering
    \includegraphics[width=0.6\textwidth]{images/graphs/fig_eda_class_imbalance.png}
    \caption{Distribution of fraudulent and non-fraudulent transactions}
    \label{fig:eda_class_imbalance}
\end{figure}


\subsection{Temporal Patterns}

Fraud occurrence exhibits clear temporal patterns, as shown in Figure~\ref{fig:eda_fraud_by_hour}. Fraud rates vary across hours of the day, with elevated rates during late-night and early-morning periods, while daytime hours show substantially lower incidence. Fraud prevalence also differs across weekdays, indicating that transaction risk is not temporally uniform. These patterns motivate the construction of temporal features and support the broader use of a sequence-based model. 

\begin{figure}[H]
    \centering
    \includegraphics[width=0.8\textwidth]{images/graphs/fig_eda_fraud_by_hour.png}
    \caption{Hourly fraud rate}
    \label{fig:eda_fraud_by_hour}
\end{figure}

\subsection{Transaction Amount}

As shown in figure~\ref{fig:eda_amount_by_class}, transaction amounts are highly right-skewed and span several orders of magnitude. Fraudulent transactions tend to involve higher amounts and greater dispersion than legitimate ones, although substantial overlap between classes remains, indicating that transaction amount alone is insufficient for reliable fraud detection.

\begin{figure}[H]
    \centering
    \includegraphics[width=0.75\textwidth]{images/graphs/fig_eda_amount_by_class.png}
    \caption{Distribution of transaction amounts by class on a logarithmic scale}
    \label{fig:eda_amount_by_class}
\end{figure}


\subsection{Feature Correlations}

Correlation analysis among numerical and engineered temporal features reveals only limited linear association with the fraud label. This suggests that fraud patterns are unlikely to be captured adequately by simple linear relationships alone and supports the use of more flexible nonlinear models.

\subsection{High-Cardinality Features}

Several categorical variables, including merchant, city, job, and ZIP code, exhibit high cardinality and long-tailed frequency distributions. As shown in Figure~\ref{fig:eda_fraud_by_category}, fraud rates vary substantially across merchant categories, indicating heterogeneous category-level risk. These properties motivate the use of encoding strategies specifically designed for high-cardinality categorical features.

\begin{figure}[H]
    \centering
    \includegraphics[width=0.9\textwidth]{images/graphs/fig_eda_fraud_by_category_topN.png}
    \caption{Fraud rates across top merchant categories}
    \label{fig:eda_fraud_by_category}
\end{figure}

\section{Feature Engineering}
\label{sec:feature_engineering}
This section describes the feature construction choices used to represent transactions in a form suitable for modelling. The goal is not to maximize the number of variables, but to retain information that may improve fraud detection while excluding fields that mainly encourage memorization or leakage.

\subsection{Temporal Feature Construction}
Transaction timestamps are transformed into interpretable temporal features. From the raw transaction datetime, we derive the hour of day, day of week, and month, together with a binary weekend indicator. These variables are intended to capture regular temporal patterns in transaction behaviour that may be relevant for fraud detection.

Because hour, weekday, and month are cyclical rather than linear quantities, they are encoded using sine and cosine transformations. This preserves the periodic structure of time by ensuring, for example, that 23:00 and 00:00 are treated as neighbouring values rather than as opposite ends of a numeric scale. For a cyclical variable \( t \) with period \( T \), the transformation is given by
\[
\sin\left( \frac{2\pi t}{T} \right), \quad
\cos\left( \frac{2\pi t}{T} \right).
\]

In addition, customer age is derived from date of birth by computing the difference between the transaction date and the recorded birth date in years. After these transformations, the original timestamp and date-of-birth fields are removed in order to reduce redundancy and avoid passing raw date information directly to the model.


\subsection{Categorical Feature Encoding}
Categorical variables are encoded according to their cardinality. Low-cardinality variables are represented using one-hot encoding, which preserves category membership without imposing an arbitrary numerical ordering. High-cardinality variables, such as merchant, city, job, and ZIP code, are encoded using target encoding estimated exclusively on the training data. This choice reduces dimensionality relative to one-hot encoding while still retaining category-level risk information. To prevent leakage, the target-encoding mappings are learned only on the training split and then applied unchanged to validation and test observations.

\subsection{Feature Exclusion and Leakage Prevention}
\label{sec:feature_exclusion}

Several identifier-like variables, including credit card numbers, transaction identifiers, names, street-level addresses, and raw timestamps, are excluded from the feature set. Although these variables may appear highly predictive within the observed sample, they primarily allow the model to memorize specific entities or historical incidents rather than learn transferable behavioural patterns.

This phenomenon is a well-documented form of target leakage and entity memorization in supervised learning. \textcite{kaufman_leakage_2012} demonstrate that high-cardinality identifiers often exhibit strong apparent predictive power due to data collection artifacts and sampling bias, leading to overly optimistic performance estimates. In the context of fraud detection, \textcite{dal_pozzolo_credit_2015} emphasize that repeated entity occurrences can bias evaluation and hinder generalization under temporal and behavioral drift.

Accordingly, identifier-based variables are removed in order to reduce leakage risk and to ensure that the model is evaluated on its ability to detect fraud-related behaviour rather than recall specific entities from the historical sample.

\section{Data Splitting Strategy}
\label{sec:data_split}

To ensure a robust evaluation, we follow a three way splitting strategy that mirrors the workflow implemented in our Python pipeline. The dataset is provided in two separate files, \texttt{fraudTrain.csv} and \texttt{fraudTest.csv}, which are temporally disjoint. We interpret \texttt{fraudTrain.csv} as the historical development sample used for model building and hyperparameter tuning, while \texttt{fraudTest.csv} is treated as a final held out evaluation sample that is never used during model selection.

\subsection{Development Sample: Training and Validation Split}

All preprocessing and model selection steps are performed exclusively on \texttt{fraudTrain.csv}. From this development dataset, we create an internal validation set via a stratified random split:
\[
\begin{aligned}
(X_{\text{train}}, y_{\text{train}}),\ (X_{\text{val}}, y_{\text{val}})
&\leftarrow \texttt{train\_test\_split}(X, y, \texttt{test\_size}=0.2, \\
&\qquad\texttt{stratify}=y,\ \texttt{random\_state}=42).
\end{aligned}
\]
Stratification preserves the severe class imbalance in both subsets, ensuring that the fraud prevalence in the validation sample remains representative of the development data. The training subset (\(80\%\)) is used to fit model parameters, including all preprocessing transformations, while the validation subset (\(20\%\)) is used for hyperparameter selection and model comparison. This results in an 80/20 split between training and validation data within the original Kaggle training set.

\subsection{External Test Sample}

The file \texttt{fraudTest.csv} is used as an external test set \((X_{\text{test}}, y_{\text{test}})\). This split is defined by the dataset provider and is temporally separated from the training data, which supports a realistic out of sample assessment under potential temporal variation. Importantly, \texttt{fraudTest.csv} is not used for hyperparameter tuning, feature selection, calibration decisions, or any other model development step. It is only used once at the end to report final model performance.

\subsection{Preprocessing within the Split and Leakage Prevention}

All feature engineering steps, including the construction of cyclical time features (sine and cosine transformations) and the derivation of age from date of birth, are applied consistently to the training, validation, and test sets. However, to avoid information leakage, model fitting related transformations are learned only on the training portion of the development sample. Concretely, the preprocessing pipeline includes:
\begin{itemize}
    \item Standardization of numerical features using \texttt{StandardScaler},
    \item One hot encoding for low cardinality categorical variables using \texttt{OneHotEncoder},
    \item Target encoding for high cardinality categorical variables using \texttt{TargetEncoder}.
\end{itemize}
These transformations are implemented within a single pipeline object, such that \texttt{fit} is applied only on \((X_{\text{train}}, y_{\text{train}})\) and the resulting transformations are subsequently transformed and applied to \(X_{\text{val}}\) and \(X_{\text{test}}\). This ensures that scaling parameters, one hot encoding categories, and target encoding mappings are estimated solely from the training subset and do not incorporate information from validation or test observations.

\subsection{Model Selection versus Final Evaluation}

The internal validation set is used for model selection. For the LSTM, hyperparameter selection focuses on the optimiser learning rate and the dropout rate, which are the two settings most directly associated with sequential learning and generalisation under class imbalance. The remaining structural hyperparameters, such as the hidden size, the number of recurrent layers, the projection dimension, the sequence length, and the batch size, are held fixed at values motivated in Section~\ref{sec:lstm}. Hyperparameter selection is based on validation PR-AUC, which is preferred over ROC-AUC and accuracy because it directly reflects ranking quality on the minority class. The exact grid and selection procedure are documented in Section~\ref{subsec:lstm_grid}.

After selecting the best configuration, we report two sets of results:
\begin{itemize}
    \item \textbf{Internal validation performance} on \((X_{\text{val}}, y_{\text{val}})\), which reflects model selection quality within the development sample,
    \item \textbf{External test performance} on \((X_{\text{test}}, y_{\text{test}})\), which reflects final generalization to unseen and temporally separated data.
\end{itemize}
This separation between development and final evaluation is critical in fraud detection settings, where class imbalance and potential temporal drift can otherwise lead to overly optimistic performance estimates if the test set is used repeatedly during tuning.

\section{Handling Class Imbalance}
\label{sec:imbalance}

Class imbalance poses a significant challenge for supervised learning algorithms, as standard classifiers tend to be biased toward the majority class. To examine which interventions, if any, are needed in the present setting, we evaluate three complementary training regimes: training on the imbalanced data with plain binary cross-entropy (no explicit imbalance handling), training with cost-sensitive binary cross-entropy that reweights the positive-class term in the loss, and training on a SMOTE-resampled version of the input sequences. The latter two are described in detail below; their empirical comparison against the no-intervention reference is reported in Section~\ref{subsec:smote}.

\subsection{Cost-sensitive Loss}
\label{subsec:cost_sensitive_loss}

The primary imbalance-handling mechanism in this study is a cost-sensitive binary cross-entropy loss. For a logit $z$ and a binary label $y \in \{0, 1\}$, the per-sample loss is
\begin{equation}
\mathcal{L}(z, y) \;=\; -\,w_{+}\, y \,\log \sigma(z) \;-\; (1 - y)\, \log\!\big(1 - \sigma(z)\big),
\end{equation}
where $\sigma(\cdot)$ denotes the logistic function and $w_{+}$ is a positive-class weight that scales the penalty incurred when a fraudulent transaction is missed. Following the standard cost-sensitive convention, the weight is set to the empirical class ratio
\begin{equation}
w_{+} \;=\; \frac{N_{\text{neg}}}{N_{\text{pos}}},
\end{equation}
which makes the contribution of the minority class to the loss approximately equal to that of the majority class. On the training subset constructed in Section~\ref{sec:data_split}, this yields $w_{+} \approx 150$. This formulation is implemented as PyTorch's \texttt{BCEWithLogitsLoss} with the \texttt{pos\_weight} argument and is the deep-learning analogue of the class-weight adjustments commonly used for tree-based and linear classifiers. Because no synthetic data is introduced, the training distribution preserves the original behavioural patterns of fraudulent and legitimate transactions, and any subsequent decision-threshold tuning operates on probability estimates that have not been distorted by oversampling.

\subsection{SMOTE}
\label{subsec:smote_method}

As an alternative to cost-sensitive training, we also evaluate the Synthetic Minority Over-sampling Technique (SMOTE), originally proposed by \textcite{chawla_smote_2002}, which generates artificial minority-class samples via interpolation in feature space.

For each minority class observation $\mathbf{x}_i \in \mathbb{R}^d$, $k$ nearest neighbors are identified within the minority class using Euclidean distance. Synthetic samples are generated by selecting a neighbor $\mathbf{x}_{\text{nn}}$ at random and computing
\begin{equation}
\mathbf{x}_{\text{synthetic}} = \mathbf{x}_i + \lambda \left( \mathbf{x}_{\text{nn}} - \mathbf{x}_i \right),
\end{equation}
where $\lambda \sim \mathcal{U}(0,1)$. This procedure places new samples along line segments between neighboring minority observations, promoting more general decision regions.

In this study, SMOTE is implemented using the \texttt{scikit-learn}-compatible implementation from the \texttt{imbalanced-learn} library, which closely follows the original formulation of \textcite{chawla_smote_2002}, with $k = 5$ by default. Because the LSTM operates on per-card sliding windows rather than on individual transactions, SMOTE is applied at the sequence level: each window is flattened into a single high-dimensional feature vector, the interpolation procedure above is run in that flattened space, and the resulting synthetic vectors are reshaped back into the original sequence layout. The minority class is up-sampled to ten percent of the majority class, which is the configuration used in the experiments reported in Section~\ref{subsec:smote}; full balancing was found to exhaust available memory on the training machine and was not used. Synthetic samples are generated only from the training portion of the development sample to avoid information leakage; validation and test data preserve the original class distribution. When SMOTE is used, the positive-class weight in the loss is set to one, so that the loss reduces to standard binary cross-entropy on a now-balanced training set and the two imbalance-handling mechanisms do not compound.

This design allows us to compare cost-sensitive and SMOTE-based LSTM models under identical evaluation conditions while preserving the original class distribution in the held-out test data.


\section{Sequence Construction}
\label{sec:sequence_construction}

After feature engineering and preprocessing, transactions are arranged into per-card sequences before being passed to the model. This step is the bridge between the row-wise representation used by the classical baselines in our previous study and the sequence-aware representation required by the LSTM. The construction procedure is deterministic and operates exclusively within each split, so that no information crosses the train, validation, or test boundary.

For every credit-card identifier \(\texttt{cc\_num}\), the corresponding transactions are sorted in ascending order of \texttt{trans\_date\_trans\_time}. For each transaction \(i\) of card \(c\), a fixed-length window is then constructed from the up to \(L = 20\) most recent transactions of the same card with a timestamp less than or equal to that of transaction \(i\), inclusive. The label associated with the window is the \texttt{is\_fraud} indicator of transaction \(i\); the labels of the preceding transactions in the window are not used as features. When fewer than \(L\) prior transactions are available for the same card, the window is left-padded with zero vectors and the original number of valid transactions is recorded as a length mask, so that the recurrent model can ignore padded positions during training and inference.

Importantly, \(\texttt{cc\_num}\) is used only as a grouping key for sequence construction and is not included in the feature set. As discussed in Section~\ref{sec:feature_exclusion}, the card identifier is excluded from \texttt{FEATURE\_COLS} for leakage-prevention reasons, and this exclusion is preserved here. The card identity therefore never enters the model tensor; what the model sees is the same row-level feature representation as the classical baselines, with the additional structure that each prediction is conditioned on the recent transaction history of the same card. This distinction allows the sequence model to capture per-card behavioural context without re-introducing the entity-memorisation risk associated with high-cardinality identifiers.

After preprocessing, each transaction is represented by 78 numerical features. Applied to the temporal split described in Section~\ref{sec:data_split}, the construction procedure yields 400{,}425 training and 70{,}664 validation sequences, of which 2{,}655 and 247 are fraudulent, respectively, and 555{,}719 test sequences with 2{,}145 fraudulent. These counts equal the row counts of the underlying splits because every transaction generates exactly one window centred on itself; the windowing step does not change the number of training instances, only the representation of each instance.


\section{LSTM}
\label{sec:lstm}

\subsection{Model Choice and Motivation}
Long Short-Term Memory (LSTM) networks are a class of recurrent neural networks designed for sequential data. In credit card fraud detection, this is relevant because suspicious activity may arise not only from the characteristics of a single transaction, but also from its position within a sequence of prior transactions. An LSTM is therefore a natural candidate when the goal is to model short-run and longer-run behavioural patterns in transaction streams.

A standard recurrent neural network passes a hidden state from one time step to the next, but it often struggles when useful information lies far back in the sequence. As the sequence becomes longer, the signal from earlier steps can fade, which makes learning difficult. This problem is commonly described as the vanishing gradient problem. \textcite{olah_understanding_2015} explain that this is one of the main reasons why standard RNNs have difficulty learning long term dependencies.

For this reason, the present study uses an LSTM model. LSTMs were introduced to preserve useful information over longer parts of a sequence by adding a separate memory path called the cell state and by using gates that control what information is kept, added, and passed forward. This makes the model well suited to the present setting, where fraud may be linked to temporal patterns and changes in behaviour across a sequence of transactions rather than to a single feature value alone.

\subsection{Model Specification}

Let $x_1, x_2, \dots, x_T$ denote a sequence of transaction feature vectors. At each time step $t$, the LSTM receives the current input $x_t$, the previous hidden state $h_{t-1}$, and the previous cell state $C_{t-1}$. The model then computes

\[
f_t = \sigma \left( W_f [h_{t-1}, x_t] + b_f \right), \quad \text{(forget gate)}
\]

\[
i_t = \sigma \left( W_i [h_{t-1}, x_t] + b_i \right), \quad \text{(input gate)}
\]

\[
\tilde{C}_t = \tanh \left( W_C [h_{t-1}, x_t] + b_C \right), \quad \text{(candidate cell state)}
\]

\[
C_t = f_t \odot C_{t-1} + i_t \odot \tilde{C}_t, \quad \text{(cell state update)}
\]

\[
o_t = \sigma \left( W_o [h_{t-1}, x_t] + b_o \right), \quad \text{(output gate)}
\]

\[
h_t = o_t \odot \tanh(C_t), \quad \text{(hidden state)}
\]

The forget gate determines which parts of the old memory are retained. The input gate determines which new information enters the memory. The output gate determines which parts of the updated memory are exposed to the hidden state. In contrast to a standard RNN, where the hidden state is fully overwritten at each step, the cell state is updated through a gated additive rule. Old information is retained or discarded by the forget gate, and new information is added through the input gate. This gives the model a more stable way to carry information through time and helps it learn longer range dependencies \parencite{olah_understanding_2015}.

In practical terms, the LSTM processes a sequence of transaction feature vectors and updates an internal representation of recent behaviour over time. This is useful in fraud detection because the model can evaluate whether the current transaction is consistent with earlier activity or constitutes a suspicious deviation. Since the transaction data are transformed into numerical feature vectors during preprocessing, the LSTM operates directly on structured transactional inputs.

\subsection{Network Architecture}
\label{subsec:lstm_arch}

The model implemented in this study consists of three components arranged in sequence. First, an input projection layer maps each 78-dimensional transaction feature vector to a 64-dimensional dense representation; this step decouples the input feature dimension from the recurrent hidden size and provides a small amount of learned linear preprocessing before the recurrent dynamics begin. Second, a stack of two LSTM layers with hidden size 128, applied in the unidirectional configuration, processes the projected sequence. The two-layer depth is sufficient to learn compositional temporal patterns over windows of length 20 without the parameter and runtime cost of a deeper stack, while still benefitting from the gating mechanisms described in Section~\ref{sec:lstm}. Third, the hidden state at the final time step is passed through a small dense classification head that applies a dropout mask, projects the 128-dimensional representation to 64 dimensions through a ReLU non-linearity, and finally maps to a single scalar logit. A sigmoid activation is applied at evaluation time to produce the fraud probability $\hat{p}(y = 1 \mid x_{1:T})$; during training, the logit is consumed directly by the loss function described below.

\subsection{Training Procedure}
\label{subsec:lstm_train}

The model is implemented in PyTorch and optimised with AdamW using a learning rate of $10^{-3}$ and weight decay of $10^{-5}$. Gradient norms are clipped at 1.0 to stabilise training. The objective is binary cross-entropy. Three configurations of the loss are evaluated and compared on the validation set in Section~\ref{subsec:smote}: plain binary cross-entropy on the imbalanced data (no positive-class weighting), cost-sensitive binary cross-entropy with the positive-class weight set to the empirical ratio of negatives to positives in the training subset (which evaluates to approximately 150) as introduced in Section~\ref{subsec:cost_sensitive_loss}, and standard binary cross-entropy on a SMOTE-rebalanced training set with the positive-class weight reset to one, as specified in Section~\ref{subsec:smote_method}. The plain binary cross-entropy configuration is selected as the baseline LSTM on the basis of the validation comparison reported in Section~\ref{subsec:smote}; the cost-sensitive and SMOTE variants are reported alongside the baseline in the test-set comparison of Section~\ref{subsec:model_comparison} for transparency.

Training is run for at most 30 epochs with a mini-batch size of 256 sequences. Early stopping monitors validation PR-AUC at the end of every epoch with a patience of five: the model state corresponding to the highest observed validation PR-AUC is retained, and training terminates once no improvement is observed for five consecutive epochs. This protocol selects the checkpoint with the strongest ranking quality on the minority class and provides an automatic safeguard against the mild overfitting visible in the training-loss trajectory, which is examined in Section~\ref{sec:training_results}.

After training and checkpoint selection, a single decision threshold is chosen on the validation set by maximising the F$_2$ score on the validation probabilities. The choice of F$_2$ over F$_1$ reflects the operational asymmetry of fraud detection, in which a missed fraudulent transaction is more costly than a false-positive review. The selected threshold is then frozen and applied unchanged to the test set, in line with the development-versus-evaluation separation described in Section~\ref{sec:data_split}. Test-set metrics are reported both at this F$_2$-tuned threshold and at the default threshold $\tau = 0.5$, so that any contribution of threshold tuning to the headline results is transparent.

\subsection{Hyperparameter Grid Search}
\label{subsec:lstm_grid}

Within the broader hyperparameter inventory of the model (number of recurrent layers, hidden size, dropout, learning rate, batch size, sequence length), two are varied empirically and the remainder are held fixed at the values described in Sections~\ref{subsec:lstm_arch} and \ref{subsec:lstm_train}. The two varied hyperparameters are the optimiser learning rate and the dropout rate applied between the recurrent layers and within the classification head: the learning rate governs the optimisation regime under the large positive-class weight, while dropout is the most direct lever on the train/validation gap that motivates regularisation in the first place. The grid search is implemented as a small full-factorial grid over
\[
\text{lr} \in \{3\times 10^{-4},\ 1\times 10^{-3},\ 3\times 10^{-3}\}
\quad \text{and} \quad
\text{dropout} \in \{0.3,\ 0.5\},
\]
yielding six trials. The lower dropout value coincides with the configuration used as the default in our implementation and serves as a reference point; the higher value increases regularisation. The learning-rate values are spaced geometrically around the default and bracket both more conservative and more aggressive optimisation regimes.

Every trial is trained from the same initialisation seed, so that observed differences across trials reflect the varied hyperparameters rather than initialisation noise. An explicit estimate of run-to-run variance is therefore beyond the scope of this experiment. Each trial is otherwise identical to the baseline configuration: the same training, validation, and test sequences, the same loss function and optimiser, the same early-stopping protocol with patience five, and the same F$_2$-based threshold-selection step on the validation set. The final configuration carried forward to the test-set comparison is the one with the highest validation PR-AUC.




\pagebreak
\chapter{Results}
\section{Training}
\label{sec:training_results}

The training procedure, network architecture, loss function, and threshold-selection protocol are specified in Section~\ref{sec:lstm}. This section reports the training dynamics observed when the baseline configuration is fit on the development sample.

Figure~\ref{fig:lstm_training_curves} reports the training dynamics of the baseline configuration. The validation PR-AUC rises rapidly within the first few epochs and reaches its maximum well before the maximum number of epochs is exhausted, after which validation performance plateaus while the training loss continues to decrease. This divergence between training and validation loss is consistent with mild overfitting and motivates the use of early stopping as well as the dropout-based regularisation explored in the hyperparameter grid search below.

\begin{figure}[H]
    \centering
    \includegraphics[width=0.9\textwidth]{images/graphs/lstm_training_curves.png}
    \caption{Training dynamics of the baseline LSTM. Loss (left) and PR-AUC (right) are reported for the training and validation sets across epochs. The dashed vertical line marks the epoch with the highest validation PR-AUC, which is the checkpoint retained for evaluation.}
    \label{fig:lstm_training_curves}
\end{figure}

\subsection{Hyperparameter Tuning}
\label{subsec:lstm_hpt}

The grid search and selection protocol are described in Section~\ref{subsec:lstm_grid}. Table~\ref{tab:lstm_hpt} reports the resulting validation PR-AUC across the six configurations. To preserve the separation between development and final evaluation, only validation results are reported here; the test set is reserved for the comparison in Section~\ref{subsec:model_comparison}.

\begin{table}[H]
\centering
\caption{Hyperparameter grid search over learning rate and dropout. Selection is based on validation PR-AUC; the seed is held constant across configurations.}
\label{tab:lstm_hpt}
\small
\renewcommand{\arraystretch}{1.2}
\begin{tabular}{lcc}
\hline
\textbf{lr} & \textbf{dropout} & \textbf{Val PR-AUC} \\
\hline
$3\times 10^{-4}$ & 0.3 & 0.830 \\
$1\times 10^{-3}$ & 0.3 & \textbf{0.889} \\
$3\times 10^{-3}$ & 0.3 & 0.878 \\
$3\times 10^{-4}$ & 0.5 & 0.828 \\
$1\times 10^{-3}$ & 0.5 & 0.874 \\
$3\times 10^{-3}$ & 0.5 & 0.888 \\
\hline
\end{tabular}
\end{table}

The configuration with the highest validation PR-AUC is the combination $\text{lr} = 1\times 10^{-3}$ and $\text{dropout} = 0.3$. Reducing the learning rate to $3\times 10^{-4}$ degrades validation PR-AUC by approximately six percentage points, indicating that the lowest-learning-rate trials are under-trained within the available epoch budget. Increasing the learning rate to $3\times 10^{-3}$ slightly worsens the validation score relative to the chosen value, and increasing dropout to $0.5$ produces a small but consistent reduction across all learning rates. The differences between the six configurations are small, and the seed is held constant across runs, so the observed variation reflects the varied hyperparameters rather than an estimate of run-to-run variability. The selected configuration is carried forward into Section~\ref{subsec:smote}, in which the loss-function variants are compared at these fixed architecture and optimisation hyperparameters.



\subsection{Effect of Class Imbalance Handling}
\label{subsec:smote}

To select the loss configuration carried forward into the test-set comparison, three variants of the LSTM are trained and evaluated on the validation split: a reference model trained with plain binary cross-entropy on the imbalanced data (no imbalance handling), the cost-sensitive variant introduced in Section~\ref{subsec:cost_sensitive_loss} that uses $\text{pos\_weight} \approx 150$, and the sequence-level SMOTE variant specified in Section~\ref{subsec:smote_method}. All three variants share the architecture, optimiser, and hyperparameters identified in Section~\ref{subsec:lstm_hpt} and differ only in the imbalance-handling mechanism applied during training. Table~\ref{tab:lstm_imbalance_val} reports the validation PR-AUC at the best epoch and the F$_2$-optimal decision threshold estimated on each variant's validation probabilities; both quantities are computed on the validation split alone and are reused unchanged at test time.

\begin{table}[H]
\centering
\caption{Validation comparison of the three imbalance-handling variants of the LSTM. Validation PR-AUC is the primary selection metric; the F$_2$-optimal threshold is estimated on each variant's validation probabilities and is reused unchanged at test time.}
\label{tab:lstm_imbalance_val}
\small
\renewcommand{\arraystretch}{1.2}
\begin{tabular}{lcc}
\hline
\textbf{Loss configuration} & \textbf{Val PR-AUC} & \textbf{F$_2$-tuned threshold} \\
\hline
Plain BCE (no imbalance handling) & \textbf{0.910} & 0.181 \\
Cost-sensitive BCE ($\text{pos\_weight} \approx 150$) & 0.889 & 0.993 \\
SMOTE-rebalanced training set & 0.852 & 0.144 \\
\hline
\end{tabular}
\end{table}

The plain-BCE configuration achieves the highest validation PR-AUC of the three, exceeding the cost-sensitive variant by roughly two percentage points and the SMOTE variant by roughly six. The F$_2$-tuned threshold differs sharply across configurations: the cost-sensitive variant requires a near-extreme threshold of $\tau \approx 0.993$ to land on a usable operating point, while the plain-BCE and SMOTE variants land on more moderate thresholds. This difference reflects the large positive-class weight in the cost-sensitive loss, which pushes logits upward and produces a strongly bimodal score distribution; the F$_2$ tuner has to compensate by selecting a threshold that almost touches one. The plain-BCE model, by contrast, retains probability calibration and accordingly tunes to a moderate threshold near $0.18$. The plain-BCE configuration is therefore selected as the baseline LSTM for the test-set comparison reported in Section~\ref{subsec:model_comparison}; the cost-sensitive and SMOTE variants are reported alongside it for transparency.



\section{Model Comparison on the Test Set}
\label{subsec:model_comparison}

Table~\ref{tab:lstm_vs_classic} reports the held-out test performance of the validation-selected baseline LSTM, alongside the two imbalance-handled LSTM variants and the Random Forest and Support Vector Machine baselines from our previous study. The held-out test file \texttt{fraudTest.csv} contains $555{,}719$ transactions ($2{,}145$ fraudulent and $553{,}574$ legitimate). Classical baselines are reported at their F$_2$-optimal validation thresholds; each LSTM variant is reported both at the default decision threshold $\tau = 0.5$ and at its F$_2$-optimal threshold tuned on the validation set, exactly as recorded in Table~\ref{tab:lstm_imbalance_val}. The validation-selected baseline LSTM is the plain binary cross-entropy variant; its rows are highlighted in bold.

\begin{table}[H]
\centering
\caption{Test-set comparison between the LSTM variants and the classical baselines. The plain-BCE LSTM, selected on validation in Section~\ref{subsec:smote}, is the baseline (in bold). The cost-sensitive and SMOTE variants are included for transparency. Classical baselines are evaluated at their F$_2$-optimal validation thresholds; each LSTM variant is reported at the default threshold and at its F$_2$-optimal validation threshold.}
\label{tab:lstm_vs_classic}
\small
\renewcommand{\arraystretch}{1.2}
\begin{tabular}{llcccccc}
\hline
\textbf{Model} & \textbf{Threshold} & \textbf{ROC-AUC} & \textbf{PR-AUC} & \textbf{Precision} & \textbf{Recall} & \textbf{F$_1$} \\
\hline
SVM           & F$_2$-tuned ($0.025$)    & 0.903 & 0.101 & 0.18 & 0.38 & 0.24 \\
Random Forest & F$_2$-tuned ($0.201$)    & 0.982 & 0.509 & 0.59 & 0.41 & 0.48 \\
\hline
\textbf{LSTM (plain BCE)} & \textbf{default ($0.500$)}    & \textbf{0.994} & \textbf{0.931} & \textbf{0.95} & \textbf{0.86} & \textbf{0.90} \\
\textbf{LSTM (plain BCE)} & \textbf{F$_2$-tuned ($0.181$)} & \textbf{0.994} & \textbf{0.931} & \textbf{0.86} & \textbf{0.90} & \textbf{0.88} \\
\hline
LSTM (cost-sensitive) & default ($0.500$)        & 0.988 & 0.899 & 0.60 & 0.91 & 0.72 \\
LSTM (cost-sensitive) & F$_2$-tuned ($0.993$)    & 0.988 & 0.899 & 0.88 & 0.85 & 0.86 \\
\hline
LSTM + SMOTE          & default ($0.500$)        & 0.986 & 0.862 & 0.77 & 0.83 & 0.80 \\
LSTM + SMOTE          & F$_2$-tuned ($0.144$)    & 0.986 & 0.862 & 0.53 & 0.89 & 0.66 \\
\hline
\end{tabular}
\end{table}

The baseline LSTM (plain BCE) dominates the classical baselines on every metric. At the F$_2$-tuned operating point, it lifts test PR-AUC from $0.509$ (Random Forest) to $0.931$, the fraud-class $F_1$ score from $0.48$ to $0.88$, precision from $0.59$ to $0.86$, and recall from $0.41$ to $0.90$. Even at the default threshold $\tau = 0.5$, the LSTM achieves $P = 0.95$, $R = 0.86$, $F_1 = 0.90$, with only $106$ false positives across the $553{,}574$ legitimate transactions in the test set. The improvement over the SVM is even larger and is consistent with the limited capacity of a linear decision boundary in this high-dimensional, highly imbalanced setting. Importantly, the LSTM dominates the Random Forest on both axes of the precision--recall trade-off simultaneously rather than buying recall at the cost of precision, as the SMOTE-trained classical models did in our previous study.

The headline gain in PR-AUC, from $0.509$ to $0.931$, is unusually large and motivates a direct interpretation of the structural difference between the model classes. The Random Forest in our previous study scored each transaction in isolation and could not condition on the recent history of the same card, because the card number itself was excluded from the feature set as a high-cardinality identifier. The LSTM does not reintroduce the card number as a feature either, but it uses the card number as a grouping key when constructing per-card sliding windows, as described in Section~\ref{sec:sequence_construction}. As a result, each prediction is conditioned on the same card's preceding transactions in time order, while the model never observes the card identity itself. This is a structural rather than a statistical advantage: the LSTM can flag a transaction as anomalous \emph{relative to the recent behaviour of the same card}, which the row-wise classical baseline is architecturally unable to do. The same property explains why the LSTM achieves both higher recall and higher precision: the additional context permits the model to suppress false positives that the row-wise baseline cannot distinguish from genuine fraud.

The two imbalance-handled LSTM variants, included in the table for transparency, also outperform the classical baselines but underperform the validation-selected baseline. The cost-sensitive variant (PR-AUC $0.899$, F$_2$-tuned $F_1 = 0.86$) loses approximately three percentage points of PR-AUC and two of $F_1$ relative to the baseline; its F$_2$-tuned threshold of $\tau \approx 0.993$ reflects the upward shift in logits introduced by the large positive-class weight, as discussed in Section~\ref{subsec:smote}. The SMOTE variant (PR-AUC $0.862$, F$_2$-tuned $F_1 = 0.66$) performs worst of the three: at its F$_2$-tuned threshold the precision drops to $0.53$, indicating that synthetic oversampling introduces a substantial false-positive burden without a corresponding gain in recall. The test results therefore replicate the validation ordering from Section~\ref{subsec:smote}.

Taken together, these findings indicate that explicit imbalance handling is not necessary to obtain strong fraud-detection performance with this LSTM, and that both standard imbalance-handling techniques measurably degrade performance in this setting. This contrasts with the classical setting reported in our previous study, in which SMOTE substantially degraded both Random Forest and SVM performance — the Random Forest's fraud-class $F_1$ fell from $0.48$ to $0.13$ and the SVM collapsed into a near-trivial classifier with $0.01$ precision — while the row-wise classifiers benefited from cost-sensitive class weights. The deep-learning setting therefore behaves differently along two dimensions. First, the per-card sequence representation introduced in Section~\ref{sec:sequence_construction} provides enough discriminative signal that the model can rank fraudulent transactions accurately even when the loss does not explicitly compensate for the class ratio of roughly $150{:}1$; in this regime, the cost-sensitive correction sacrifices probability calibration for marginal recall gains and ends up worse on $F_1$ overall. Second, sequence-level SMOTE interpolates between flattened multivariate windows, and the resulting synthetic windows may not lie on the manifold of plausible per-card transaction histories. The model cannot use these synthetic windows to learn the temporal structure that drives the headline gain reported above, and ranking quality measurably suffers. The conclusion from both studies is therefore that synthetic oversampling does not help under realistic temporal evaluation, and the additional conclusion specific to the deep-learning setting is that even cost-sensitive class weighting can be omitted when the model architecture itself captures the relevant structure.


\subsection{Statistical Significance on Fraud Recall}
\label{statistical_significance}


\pagebreak
\chapter{Discussion \& Conclusion}

\section{Discussion}



\subsection{Interpretation of Results and Comparison with Prior Work}



\subsection{Limitations}



\section{Conclusion}



%%%%%%%%%%%%%%%%%%%%%%%%%%%%%%%%%%%%%%%%%%%%%%%%%%%%%%%%%%%%%%%%%
% Ab hier werden die Verzeichnisse eingehängt
\cleardoublepage
\pagenumbering{roman}
\setcounter{page}{1}

\printbibliography
\addcontentsline{toc}{chapter}{\bibname} % Nutzt automatisch "References" aus babel


\chapter*{List of Tools}
\addcontentsline{toc}{chapter}{List of Tools}
\begin{description}
  \item[OpenAI (2025/2026).] ChatGPT (Version 5.2). \url{https://chat.openai.com/chat} \\
  – Support with programming \\
  – Generation of text suggestions

  \item[Anthropic (2025/2026).] Claude AI (Claude 3.Opus). \url{https://www.anthropic.com/claude} \\
  – Support with programming and code structuring

  \item[DeepL (2025/2026).] DeepL Translator (DeepL Pro). \url{https://www.deepl.com/translator} \\
  – Translation of individual text passages
\end{description}


%\lstlistoflistings
%\addcontentsline{toc}{section}{\lstlistlistingname} % Falls "List of Listings" korrekt ist
%%%%%%%%%%%%%%%%%%%%%%%%%%%%%%%%%%%%%%%%%%%%%%%%%%%%%%%%%%%%%%%%%
% Ab hier der Anhang eingehängt
\appendix
\chapter{Appendix}



\section{Code \& Data}
A ZIP file is provided containing the complete codebase together with the corresponding CSV data files.

% Statement of Authorship (chapter-level, bold, unnumbered)
\chapter{Statement of Authorship}
\addcontentsline{toc}{chapter}{{Statement of Authorship}}
\vspace{1cm}
"We hereby declare that this research project is our own work, that it has been created by us without the help of others, using only the sources referenced, and that we will not supply any copies of this paper to any third parties without written permission by the supervisor of this project"

At the same time, all rights to this project are hereby assigned to ZHAW Zurich University of Applied Sciences, except for the right to be identified as its author.

\vspace*{1cm}

% Place and date
\noindent
\textbf{Place, Date:} \\[1em]
Winterthur, 22.01.2026

\vspace{1cm}

% Signatures heading
\noindent
\textbf{Signatures:}

\vspace{1cm}

% Signature blocks
\noindent
\begin{minipage}[t]{0.30\textwidth}
\centering
\includegraphics[width=0.8\linewidth]{images/Signatures/Signature Francis.png} \\[0.1em]
\dotfill \\[0.1em]
Francis Du
\end{minipage}
\hfill
\begin{minipage}[t]{0.30\textwidth}
\centering
\includegraphics[width=0.8\linewidth]{images/Signatures/Signature Dominik.png} \\[0.1em]
\dotfill \\[0.1em]
Dominik Kaufmann
\end{minipage}
\hfill
\begin{minipage}[t]{0.30\textwidth}
\centering
\includegraphics[width=0.8\linewidth]{images/Signatures/Signature Max.png} \\[0.1em]
\dotfill \\[0.1em]
Max Heinzer
\end{minipage}



\end{document}
